% Paper 2, §3 — Scaling laws
% Anchors: kb/notes/scaling/{kaplan-laws,chinchilla,scaling-frontier,mu-transfer}.md
%          kb/excerpts/{kaplan2020,hoffmann2022-chinchilla,yang2022-mup,
%                       scaling-laws-precision-2024,meta-llama3,meta-llama4,
%                       deepseek-v3-training,qwen3,gemini2-5}.md

\section{Scaling laws}
\label{sec:scaling-laws}

A \glsterm{scaling-law}{scaling law} is a fitted functional form $L(N, D, C)$ relating
\glsterm{cross-entropy-loss}{cross-entropy loss} to parameter count, training tokens, and training
compute, together with a closed-form recipe for how to spend a fixed
$C$ across $N$ and $D$. The fitted form is empirical; the recipe is
prescriptive; the two are routinely conflated, and most of the
``Kaplan-vs-Chinchilla'' confusion in 2022--2024 was conflation. The
thesis of this section is that the \emph{functional forms} from
\citet{kaplan2020} and \citet{hoffmann2022-chinchilla} are essentially
agreed; the \emph{compute-allocation prescription} flipped from
``train very large models, stop well short of convergence''
\citep[\S 1.1]{kaplan2020} to ``scale \glsterm{parameters}{parameters} and tokens equally''
\citep[\S 3.4]{hoffmann2022-chinchilla} on the strength of one
methodological correction (matching the LR-schedule horizon to the
actual training-token count); and the \emph{frontier-2026 practice}
sits an order of magnitude beyond Chinchilla on tokens-per-parameter
because inference cost has displaced pre-training cost as the dominant
objective \citep[\S 3.2]{deepseek-v3}. We then layer on
\textmu-Transfer \citep{yang2022-mup} as the hyperparameter-transfer
technology that makes any of these recipes tunable at frontier scale,
and the precision-aware extension of
\citet{scaling-laws-precision-2024} as the open precision $\times$ data
trade-off that the over-training strategy collides with.

The compute approximation we will use throughout is the
\citet[\S 2.1, eq.\ 2.1]{kaplan2020}\footnote{KB excerpt:
\texttt{kb/excerpts/kaplan2020.md\#sec-2-1-definitions}.} per-token
forward FLOP count $C_{\text{forward}} \approx 2N$, with the backward
pass roughly twice the forward, summing to
\begin{equation}
C \;\approx\; 6\,N\,D \quad \text{FLOPs}
\label{eq:6nd}
\end{equation}
total non-embedding training compute, valid when $d_{\text{model}}
\gtrsim n_{\text{ctx}}/12$ \citep[\S 2.1]{kaplan2020}. Every scaling
result in this section is stated against \cref{eq:6nd}.

\subsection{Kaplan power laws}
\label{sec:scaling-kaplan}

\citet{kaplan2020} fitted \glsterm{decoder-only}{decoder-only} \glsterm{transformer}{Transformer} LMs on WebText2
across $N \in [768,\, 1.5 \times 10^9]$ \glsterm{non-embedding-parameters}{non-embedding parameters} and
$D \in [22\text{M},\, 23\text{B}]$ tokens, GPT-2 \glsterm{bpe}{BPE} tokenizer,
context length $1024$, Adam (Adafactor at $N > 1\text{B}$), and a
cosine LR schedule with $3000$-step \glsterm{warmup}{warmup} decaying to zero
\citep[\S 2.2]{kaplan2020}. The headline finding is that
\glsterm{cross-entropy-loss}{cross-entropy loss} obeys a \emph{power law} in each of $N, D, C_{\min}$
when the other two are not the bottleneck.

\paragraph{The three single-axis laws.}
Verbatim from \citet[\S 1.2, eqs.\ 1.1--1.3]{kaplan2020}\footnote{KB
excerpt: \texttt{kb/excerpts/kaplan2020.md\#sec-1-2-equations}.}:
\begin{align}
L(N) &\;=\; \left(\frac{N_c}{N}\right)^{\alpha_N},
   & \alpha_N \approx 0.076,\;\; N_c \approx 8.8 \times 10^{13}, \tag{1.1} \\
L(D) &\;=\; \left(\frac{D_c}{D}\right)^{\alpha_D},
   & \alpha_D \approx 0.095,\;\; D_c \approx 5.4 \times 10^{13}, \tag{1.2} \\
L(C_{\min}) &\;=\; \left(\frac{C_c^{\min}}{C_{\min}}\right)^{\alpha_C^{\min}},
   & \alpha_C^{\min} \approx 0.050,\;\; C_c^{\min} \approx 3.1 \times 10^{8}\,\text{PF-days}. \tag{1.3}
\end{align}
Here $N$ counts \emph{non-embedding} \glsterm{parameters}{parameters}
\citep[\S 1.3]{kaplan2020},\footnote{KB excerpt:
\texttt{kb/excerpts/kaplan2020.md\#sec-1-3-notation}.}
$D$ is tokens, $C_{\min}$ is the compute that would be expended at a
batch much smaller than the critical-batch power-law $B_{\text{crit}}(L)$
\citep[\S 5.1]{kaplan2020}. The exponents are small: $\alpha_N \approx
0.076$ means doubling $N$ at fixed $D$ improves loss by a factor
$2^{-0.076} \approx 0.949$, i.e.\ $\approx 5\%$ relative.

\paragraph{The joint $L(N, D)$ form.}
The two single-axis laws are unified into a single equation
\citep[\S 1.2, eq.\ 1.5]{kaplan2020}\footnote{KB excerpt:
\texttt{kb/excerpts/kaplan2020.md\#sec-1-2-joint}.}:
\begin{equation}
L(N, D) \;=\; \left[\left(\frac{N_c}{N}\right)^{\alpha_N/\alpha_D} + \frac{D_c}{D}\right]^{\alpha_D}.
\label{eq:kaplan-joint}
\tag{1.5}
\end{equation}
The exponent combination $\alpha_N/\alpha_D \approx 0.74$ is what
delivers Kaplan's sublinear-in-$N$ data prescription: to avoid the
overfitting penalty, $D$ should grow as $D \propto N^{0.74}$
\citep[\S 1.2]{kaplan2020}. Doubling $N$ requires only $\sim 1.67\times$
more data.

\paragraph{The compute-allocation rule.}
Combined with the critical-batch power law and \cref{eq:6nd},
\cref{eq:kaplan-joint} implies a closed-form for how a fixed compute
budget $C_{\min}$ should be split across $N$, batch size $B$, and
training steps $S$ \citep[\S 1.2, eqs.\ 1.7--1.8]{kaplan2020}:
\begin{equation}
N \propto C_{\min}^{\alpha_C^{\min}/\alpha_N}, \quad
B \propto C_{\min}^{\alpha_C^{\min}/\alpha_B}, \quad
S \propto C_{\min}^{\alpha_C^{\min}/\alpha_S}, \quad
D = B \cdot S.
\label{eq:kaplan-alloc}
\tag{1.7}
\end{equation}
Plugging in the measured exponents \citep[\S 1.2]{kaplan2020}:
\begin{equation}
N \propto C_{\min}^{0.73}, \quad
B \propto C_{\min}^{0.24}, \quad
S \propto C_{\min}^{0.03}, \quad
D \propto C_{\min}^{0.27}.
\label{eq:kaplan-073}
\end{equation}
Reading: a $10\times$ compute budget should buy a $\sim 5.4\times$
larger model, $\sim 1.7\times$ larger batch, and only $\sim 1.07\times$
more steps --- almost no extra training time, dramatically more
\glsterm{parameters}{parameters}. This is the prescription the GPT-3 generation took
literally: $175\text{B}$ \glsterm{parameters}{parameters} trained on $\sim 300\text{B}$
tokens, tokens-per-parameter $\approx 1.7$. Findings 1--3 of
\citet[\S 1.1]{kaplan2020}\footnote{KB excerpt:
\texttt{kb/excerpts/kaplan2020.md\#sec-1-1-bullets}.} (performance
depends on scale not shape; smooth power laws over six orders of
magnitude; universality of overfitting) are still considered correct.
Finding~5 (``train very large, stop short of convergence'') is the one
that flipped.

\subsection{Chinchilla and the parametric loss}
\label{sec:scaling-chinchilla}

\citet{hoffmann2022-chinchilla} re-ran the scaling sweep with one
methodological change: for each parameter count $N$, they trained $4$
models with cosine LR schedules whose horizon \emph{matched the actual
training-token count}, varying that horizon by a factor of $16$
\citep[\S 3.1]{hoffmann2022-chinchilla}.\footnote{KB excerpt:
\texttt{kb/excerpts/hoffmann2022-chinchilla.md\#sec-3-1}.} They then
attacked the compute-optimal allocation problem
\begin{equation}
\min_{N, D} L(N, D) \quad \text{subject to} \quad 6 N D = C
\label{eq:cc-problem}
\end{equation}
three independent ways and found the answers agreed within bootstrap
error.

\paragraph{Approach 1 --- minimum over training curves.}
For each $N \in \{75\text{M}, 250\text{M}, 500\text{M}, 1\text{B},
2.5\text{B}, 5\text{B}, 10\text{B}\}$, fit $4$ runs at $4$ different
horizons; at $1500$ log-spaced FLOP values, find the $(N, D)$ pair
with lowest loss. Fit power laws $N_{opt} \propto C^a$,
$D_{opt} \propto C^b$, obtaining $a = 0.50$, $b = 0.50$
\citep[\S 3.1]{hoffmann2022-chinchilla}.

\paragraph{Approach 2 --- IsoFLOP profiles.}
At $9$ fixed FLOP budgets $C \in [6 \times 10^{18},\, 3 \times 10^{21}]$,
sweep $N$ and read off the loss-vs-$N$ parabola minimum to get
$N_{opt}(C)$. Fit $a = 0.49$, $b = 0.51$
\citep[\S 3.2]{hoffmann2022-chinchilla}.\footnote{KB excerpt:
\texttt{kb/excerpts/hoffmann2022-chinchilla.md\#sec-3-2}.}

\paragraph{Approach 3 --- parametric loss fit.}
Pose a closed-form risk decomposition for the loss
\citep[\S 3.3, eq.\ 2]{hoffmann2022-chinchilla}\footnote{KB excerpt:
\texttt{kb/excerpts/hoffmann2022-chinchilla.md\#sec-3-3}.}:
\begin{equation}
\hat{L}(N, D) \;\triangleq\; E + \frac{A}{N^{\alpha}} + \frac{B}{D^{\beta}},
\label{eq:cc-paraloss}
\tag{2}
\end{equation}
verbatim from the paper. The first term $E$ is the irreducible loss
(entropy of natural text on the eval distribution); $A/N^{\alpha}$ is
the capacity-limited residual of a perfectly-trained $N$-parameter
model; $B/D^{\beta}$ is the optimization-limited residual of finite
training steps over a finite sample
\citep[\S 3.3]{hoffmann2022-chinchilla}. Fit $(A, B, E, \alpha, \beta)$
by L-BFGS on the Huber loss with $\delta = 10^{-3}$
\citep[\S 3.3, eq.\ 3]{hoffmann2022-chinchilla}:
\begin{equation}
\min_{A,B,E,\alpha,\beta} \;\sum_{\text{Runs } i}
   \mathrm{Huber}_\delta\!\left(\log \hat{L}(N_i, D_i) - \log L_i\right).
\label{eq:cc-fit}
\tag{3}
\end{equation}
Minimizing \cref{eq:cc-paraloss} under the $C \approx 6 N D$
constraint yields the compute-optimal frontier in closed form
\citep[\S 3.3, eq.\ 4]{hoffmann2022-chinchilla}:
\begin{equation}
N_{opt}(C) \;=\; G \!\left(\frac{C}{6}\right)^{a}, \qquad
D_{opt}(C) \;=\; G^{-1}\!\left(\frac{C}{6}\right)^{b},
\label{eq:cc-frontier}
\tag{4}
\end{equation}
with
\begin{equation}
G \;=\; \left(\frac{\alpha A}{\beta B}\right)^{\!1/(\alpha + \beta)},
\quad a \;=\; \frac{\beta}{\alpha + \beta}, \quad
b \;=\; \frac{\alpha}{\alpha + \beta},
\qquad a + b \;=\; 1\ \text{identically.}
\label{eq:cc-abG}
\end{equation}
The split $(a, b)$ is the empirical question. Approach 3 gives
$a = 0.46$, $b = 0.54$ \citep[\S 3.3]{hoffmann2022-chinchilla}.

\paragraph{All three approaches agree, and disagree with Kaplan.}
Table~\ref{tab:cc-table-2} summarizes
\citep[\S 3, Table 2]{hoffmann2022-chinchilla}:\footnote{KB excerpt:
\texttt{kb/excerpts/hoffmann2022-chinchilla.md\#sec-3-table-2}.}

\begin{table}[h]
\centering
\small
\begin{tabular}{lcc}
\toprule
Approach & $a$ ($N_{opt} \propto C^{a}$) & $b$ ($D_{opt} \propto C^{b}$) \\
\midrule
1.\ Minimum over training curves & $0.50$ \tiny{$(0.488, 0.502)$} & $0.50$ \tiny{$(0.501, 0.512)$} \\
2.\ IsoFLOP profiles             & $0.49$ \tiny{$(0.462, 0.534)$} & $0.51$ \tiny{$(0.483, 0.529)$} \\
3.\ Parametric $\hat{L}$ fit      & $0.46$ \tiny{$(0.454, 0.455)$} & $0.54$ \tiny{$(0.542, 0.543)$} \\
\midrule
\textbf{\citet{kaplan2020}}      & $\mathbf{0.73}$ & $\mathbf{0.27}$ \\
\bottomrule
\end{tabular}
\caption{Three independent estimates of the compute-allocation
exponents in $N_{opt} \propto C^{a}$, $D_{opt} \propto C^{b}$, all
within bootstrap of $a \approx b \approx 1/2$, vs.\ Kaplan's
$0.73/0.27$ split. Bracketed numbers are 10th/90th percentile
bootstrap intervals. After
\citet[\S 3, Table 2]{hoffmann2022-chinchilla}.}
\label{tab:cc-table-2}
\end{table}

\paragraph{The 70B vs.\ 280B test.}
\citet{hoffmann2022-chinchilla} fixed total compute at the Gopher
budget ($\approx 5.76 \times 10^{23}$ FLOPs) and trained Chinchilla at
$N = 70\text{B}$, $D = 1.4\text{T}$ (tokens-per-parameter $= 20$),
versus Gopher at $N = 280\text{B}$, $D = 300\text{B}$ (tokens-per-
parameter $\approx 1.07$) \citep[\S 4.1, Table 4]{hoffmann2022-chinchilla}.
Same depth ($L = 80$); half the width ($d_{\text{model}}: 16384 \to
8192$); doubled max LR; halved batch. Chinchilla beats Gopher on
every Pile subset, +0.6\,pp average across BIG-bench, and lifts \glsterm{mmlu}{MMLU}
from 60.0\% to 67.5\% \citep[\S 4.2]{hoffmann2022-chinchilla}.\footnote{KB
excerpt: \texttt{kb/excerpts/hoffmann2022-chinchilla.md\#sec-4-2}.}
The headline rule of thumb falls out of \cref{eq:cc-frontier}'s Approach-1
projection: tokens-per-parameter $\approx 20$ at compute optimum, with
a slow drift toward more data at very large scale --- $20.0$ at $400$M,
$20.2$ at $1$B, $21.1$ at $175$B, $21.2$ at $1$T \citep[\S 3, Table
3]{hoffmann2022-chinchilla}.\footnote{KB excerpt:
\texttt{kb/excerpts/hoffmann2022-chinchilla.md\#sec-3-table-3}.}

\subsection{Why Kaplan over-allocated to parameters --- the LR-horizon bug}
\label{sec:scaling-lrhorizon}

The methodological core of \citeauthor{hoffmann2022-chinchilla}'s
correction is one sentence \citep[\S 2]{hoffmann2022-chinchilla}\footnote{KB
excerpt: \texttt{kb/excerpts/hoffmann2022-chinchilla.md\#sec-2-disagree-kaplan}.}:

\begin{quote}
For a fixed learning rate \glsterm{cosine-schedule}{cosine schedule} to 130B tokens, the
intermediate loss estimates (for $D' \ll 130\text{B}$) are therefore
overestimates of the loss of a model trained with a schedule length
matching $D'$.
\end{quote}

In \citet{kaplan2020}'s sweep, every model was trained with a cosine
LR decay reaching zero at a fixed total token horizon, regardless of
how many tokens that individual run actually consumed
\citep[\S 2.2]{kaplan2020}. The intermediate-loss readouts at $D'
\ll 130\text{B}$ were therefore taken from a model whose LR was
\emph{still high} --- mid-training, not at its endpoint. A model
genuinely trained to $D'$ with an LR schedule that decays to zero at
$D'$ achieves a lower final loss than the same model trained to $D'$
with a $130\text{B}$-horizon schedule. Kaplan inferred his $L(D)$
scaling from the latter, biased upward, runs.

\begin{intuition}
Kaplan's setup was, in effect, ``take a partially-trained model and
ask how much loss is left.'' That penalizes smaller-$D$ runs because
they are not yet at their LR-decay endpoint --- their loss is the
loss of an in-progress run, not of a converged-to-its-token-budget
run. The lesson is that what we measure is the loss \emph{at training
endpoint}, where endpoint $\equiv$ end of LR schedule. Conflating
``tokens consumed'' with ``training endpoint'' was the trap. Returning
to canonical form: when \citet{hoffmann2022-chinchilla} matched the
LR-decay horizon to $D'$ for every run in
\cref{eq:cc-fit}, the apparent advantage of more-$N$-less-$D$ in
\cref{eq:kaplan-073} disappeared and the exponents $(a, b)$ flipped
from $(0.73, 0.27)$ to $\approx (0.50, 0.50)$.
\end{intuition}

The functional forms (1.1)--(1.5) are unaffected --- they fit either
sweep --- but the \emph{numerical exponents} extracted from
\cref{eq:kaplan-alloc} change qualitatively. This is the cleanest
example we have in \glsterm{llm}{LLM} scaling of an apparent recipe being driven by
the optimization protocol, not by the architecture. Notation note:
Kaplan's $N$ is \glsterm{non-embedding-parameters}{non-embedding parameters}
\citep[\S 1.3]{kaplan2020}; Hoffmann's $N$ is total \glsterm{parameters}{parameters}; the
constant offset is small at frontier scale (vocab embedding $\approx$
$0.5\%$--$5\%$ of total) but matters when comparing exponents and
constants across papers.

\subsection{Frontier 2024--2026 --- past Chinchilla}
\label{sec:scaling-frontier}

The post-Chinchilla industry move is to train substantially smaller
models on substantially more tokens than \cref{eq:cc-frontier}
prescribes. Llama~3 8B was trained on $\approx 15\text{T}$ tokens
\citep{meta-llama3}, giving tokens-per-parameter $\approx 1875$ ---
roughly $94\times$ over the Chinchilla optimum (\cref{tab:frontier-NDC}).
Llama~3 70B at the same $\approx 15\text{T}$-token budget is at
$\approx 214$ tokens-per-parameter, $\sim 11\times$ over Chinchilla.
\glsterm{precision-pinning}{DeepSeek-V3} trained $671\text{B}$ total / $37\text{B}$ active-MoE on
$14.8\text{T}$ tokens \citep[\S 1, \S 2]{deepseek-v3}, putting active
tokens-per-parameter at $\approx 400$.

\begin{table}[h]
\centering
\small
\begin{tabular}{lrrrr}
\toprule
Model & $N_{\text{active}}$ & $D$ (tokens) & $D / N_{\text{active}}$ & vs.\ Chinchilla \\
\midrule
Chinchilla 70B (2022)        & $70$B    & $1.4$T   & $20.0$  & $1.0\times$  \\
Llama 2 70B (2023)           & $70$B    & $2$T     & $28.6$  & $1.4\times$  \\
Llama 3 70B (2024)           & $70$B    & $\sim 15$T & $\sim 214$ & $\sim 11\times$ \\
Llama 3 8B  (2024)           & $8$B     & $\sim 15$T & $\sim 1875$ & $\sim 94\times$ \\
\glsterm{precision-pinning}{DeepSeek-V3} (2024) MoE       & $37$B    & $14.8$T  & $\sim 400$ & $\sim 20\times$ \\
\bottomrule
\end{tabular}
\caption{Tokens-per-active-parameter for selected frontier-scale
models, all post-Chinchilla. Llama 3 numbers from
\citet{meta-llama3}; DeepSeek-V3 from \citet[\S 2]{deepseek-v3}; the
Chinchilla baseline from \citet[\S 4.1]{hoffmann2022-chinchilla}. The
``vs.\ Chinchilla'' column is $D / (20\,N_{\text{active}})$.}
\label{tab:frontier-NDC}
\end{table}

\paragraph{Why over-train.}
The objective Chinchilla optimizes is pre-training compute,
$C_{\text{train}}$. The cost a deployed-model lab actually pays is
\begin{equation}
\text{Total cost} \;=\; C_{\text{train}} \;+\; N_{\text{queries-served}} \cdot C_{\text{inf-per-query}}.
\label{eq:total-cost}
\end{equation}
When $N_{\text{queries-served}}$ is large --- consumer chat products,
API serving --- the second term dominates, and the optimum shifts
toward smaller $N$ (lower $C_{\text{inf-per-query}}$) with proportionally
more $D$. Llama 3 8B's $\sim 1875$ tokens-per-parameter is the endpoint
of this re-objective for ``small dense model destined for billions of
queries.'' Crucially, the \glsterm{chinchilla-parametric-loss}{Chinchilla parametric loss} \cref{eq:cc-paraloss}
\emph{still applies} along this trajectory: $\hat{L}(N, D)$ keeps
decreasing in $D$ at fixed $N$, just sub-optimally on the
$C_{\text{train}}$ margin. Over-training is not a refutation of
Chinchilla; it is a different objective.

\paragraph{Reasoning models as a second compute axis.}
DeepSeek-R1 \citep{deepseek-r1}, the OpenAI o-series, and Gemini~2.5
\citep{gemini2-5} demonstrated that \glsterm{test-time-compute}{test-time compute} is now a
deployable scaling axis: production reasoning models expose
configurable thinking budgets (Gemini's \texttt{thinkingBudget}
parameter; Qwen 3's thinking-mode toggle via system prompt
\citep{qwen3}), running $5\!\times\!\!-\!\!100\!\times$ the
non-thinking cost per \glsterm{query}{query} for hard math, code, and agent loops. The
empirical surface for the joint $(C_{\text{train}}, C_{\text{test}})$
trade-off is documented but no closed-form \glsterm{scaling-law}{scaling law} captures it
across the multi-axis frontier as of 2026.

\subsection{\textmu P and \textmu Transfer --- HP transfer across width}
\label{sec:scaling-mup}

A scaling recipe is only as useful as the hyperparameters fed to it.
At frontier scale every full pretraining run costs millions of dollars,
so iterating an HP sweep on the target model is intractable. The
Maximal Update Parametrization (\textmu P) of \citet{yang2022-mup} is
the parametrization in which optimal hyperparameters are
\emph{width-invariant}, enabling \textmu Transfer: tune the master LR,
Adam betas, schedule shape, and per-layer multipliers on a small proxy
model, then copy them zero-shot to the full-scale target. The headline
demonstration: tuning at $40\text{M}$ proxy scale and transferring to
GPT-3 6.7B with total tuning cost equal to $7\%$ of one pretraining
run \citep[abstract]{yang2022-mup}.\footnote{KB excerpt:
\texttt{kb/excerpts/yang2022-mup.md\#abstract}.}

\paragraph{The \textmu P scaling rules.}
Each linear layer falls into one of three categories by width
behavior: \emph{input weights and biases} (finite input dim, infinite
output dim --- e.g.\ embeddings); \emph{output weights} (infinite
input dim, finite output dim --- e.g.\ the unembedding head); and
\emph{hidden weights} (infinite both --- e.g.\ attention QKV/O,
\glsterm{ffn}{FFN}). \textmu P is exactly the per-category rescaling
\citep[\S 4, Table 3]{yang2022-mup}\footnote{KB excerpt:
\texttt{kb/excerpts/yang2022-mup.md\#sec-4-table-3}.}:

\begin{table}[h]
\centering
\small
\begin{tabular}{lccc}
\toprule
                  & Input weights \& biases & Output weights         & Hidden weights \\
\midrule
Init.\ Var.       & $1/\text{fan\_in}$       & $\mathbf{1/\text{fan\_in}^{2}}$ & $1/\text{fan\_in}$ \\
SGD LR            & $\mathbf{\text{fan\_out}}$ & $\mathbf{1/\text{fan\_in}}$    & $1$ \\
Adam LR           & $1$                       & $\mathbf{1/\text{fan\_in}}$    & $\mathbf{1/\text{fan\_in}}$ \\
\bottomrule
\end{tabular}
\caption{The \textmu P scaling rules. Bold cells are the differences
from standard parametrization (SP), in which init variance is
$1/\text{fan\_in}$ everywhere and the LR is uniform. After
\citet[\S 4, Table 3]{yang2022-mup}.}
\label{tab:mup-table-3}
\end{table}

The effective per-weight LR is $\eta_W = \eta \cdot
(\text{Table~\ref{tab:mup-table-3} entry for } W)$ at width $n$, where
$\eta$ is the dimensionless master LR transferred from the proxy.
For Transformers, \textmu P adds one structural change
\citep[\S 4, Definition 4.1]{yang2022-mup}\footnote{KB excerpt:
\texttt{kb/excerpts/yang2022-mup.md\#sec-4-definition-4-1}.}:

\begin{definition}[\textmu P for \glsterm{transformer}{Transformer}; \citealp{yang2022-mup}]
The Maximal Update Parametrization (\textmu P) for a \glsterm{transformer}{Transformer} is
given by Table~\ref{tab:mup-table-3} and $1/d$ attention instead of
$1/\sqrt{d}$, i.e.\ the attention logit is calculated as $q^\top k / d$
instead of $q^\top k / \sqrt{d}$ where \glsterm{query}{query} $q$ and \glsterm{key}{key} $k$ have
dimension $d$.
\end{definition}

\begin{analogy}
\textmu P is the ``right'' parametrization in the same sense the
$1/\sqrt{n}$ factor in the CLT is the right scaling: $1/\sqrt{n}$ is
the unique rescaling for which $\frac{1}{\sqrt{n}}(x_1 + \cdots + x_n)
\to \mathcal{N}(0, 1)$ has a non-trivial limit
\citep[\S 2]{yang2022-mup}.\footnote{KB excerpt:
\texttt{kb/excerpts/yang2022-mup.md\#sec-2-clt}.} Width plays the role
of $n$. The \textmu P assignment of $(\text{init}, \text{LR},
\text{multiplier})$ per layer category is the unique stable choice
that preserves $\Theta(1)$-magnitude weight updates --- ``\glsterm{feature}{feature}
learning'' --- in the infinite-width limit; \glsterm{sequence-parallel}{SP} either has trivial
NTK-regime dynamics or blows up. Returning to math: under \glsterm{sequence-parallel}{SP}, after
even one Adam step, attention pre-\glsterm{softmax}{softmax} scores $q^\top k$ acquire
correlation across $i$, so $\mathrm{Var}(q^\top k) = \Theta(d)$ by
Law-of-Large-Numbers (not CLT), and the Vaswani $1/\sqrt{d_k}$
rescaling makes the \glsterm{logits}{logits} $\Theta(\sqrt{d})$ --- they blow up with
width \citep[\S 5]{yang2022-mup}. \textmu P's $1/d$ scaling absorbs the
extra factor exactly. Vaswani's $1/\sqrt{d_k}$ remains correct
\emph{at initialization} (uncorrelated $q, k$); \textmu P's $1/d$ is
correct \emph{during training}.
\end{analogy}

\paragraph{What transfers.}
\citet[\S 6]{yang2022-mup}\footnote{KB excerpt:
\texttt{kb/excerpts/yang2022-mup.md\#sec-6-what-transfers}.} report
that across width, master LR, Adam betas, LR schedule shape, per-layer
init multipliers, and most parameter multipliers transfer; across batch
size, sequence length, and training-step count they transfer
empirically within ``reasonable ranges'' (batch $\geq 32$, seqlen
$\geq 128$, steps $\geq 5000$). What does \emph{not} transfer:
regularization (weight decay, dropout) and any HP that interacts with
the regularization-vs-optimization trade-off, since regularization
strength depends on data size, not just on width
\citep[\S 6]{yang2022-mup}. Across depth the transfer is partial:
optimal init standard deviation does not transfer well across depth;
practical recipe is to fix init at the target depth and sweep the
remaining HPs \citep[\S 6.1]{yang2022-mup}.

\begin{contradiction}
Strict \citep[Definition 4.1]{yang2022-mup} \textmu P is rare in
production. Cerebras has explicit practitioner guides and Microsoft
(Phi-4) and Graphcore (u-\textmu P, \citealp{u-mup2024}) disclose
usage; Anthropic and OpenAI do not. Meta's Llama 4 tech report
\citep{meta-llama4} cites a ``\glsterm{metap}{MetaP}'' scheme for cross-scale HP
transfer with thin published detail. The community's working
assumption is that frontier labs apply some subset of
Table~\ref{tab:mup-table-3}'s rules (typically the LR scalings) but
keep Vaswani's $1/\sqrt{d_k}$ attention rather than \textmu P's $1/d$,
making them ``\textmu P-inspired'' rather than strict-Definition-4.1
variants. Independent positive results outside the
Microsoft/Cerebras/Graphcore ecosystem at $10\text{B}{+}$ scale are
scarce; whether the loose variants retain the unique-\glsterm{feature}{feature}-learning
guarantee is open.
\end{contradiction}

\subsection{Precision-aware scaling}
\label{sec:scaling-precision}

\citet{scaling-laws-precision-2024} extend the Chinchilla parametric
form \cref{eq:cc-paraloss} with precision as a third axis. The core
construction reduces to an \emph{effective parameter count}
$N_{\text{eff}}(N, b_{\text{train}})$ that is decreasing in $N$ when
training precision $b_{\text{train}}$ falls below a threshold:
training in low precision is, in this picture, equivalent to training
a smaller-$N$ model in full precision, with the equivalence-mapping
fitted from sub-frontier sweeps. Combined with a post-training
quantization-loss model, the joint result is that \emph{over-trained
models are harder to quantize} --- the more $D$ a model has consumed
relative to $N$, the larger the post-quantization loss penalty at a
given deployment precision.\deepencite{scaling-laws-precision-2024 \S
4: effective-parameter form $N_{\text{eff}}(N, b_{\text{train}})$ and
the over-training $\to$ harder-to-quantize prediction.}

This collides directly with the over-training strategy of
\cref{tab:frontier-NDC}: the same Llama-3 8B that gains inference-cost
efficiency from being small dense + token-heavy is, by the
\citet{scaling-laws-precision-2024} prediction, more difficult to
quantize at deployment time --- partially offsetting the very gain
that motivated the strategy. The industry response has been native
low-precision training: pre-train at the precision you will deploy,
not above it. \glsterm{precision-pinning}{DeepSeek-V3} reports zero irrecoverable loss spikes over
$14.8\text{T}$ tokens of \glsterm{fp8}{FP8} pre-training
\citep[abstract]{deepseek-v3}; Llama 4 likewise discloses \glsterm{fp8}{FP8}
pre-training. u-\textmu P \citep{u-mup2024} integrates \textmu P with
unit-scaling so per-tensor activation magnitudes are predictable,
enabling \glsterm{fp8}{FP8} without dynamic rescaling. The full multi-axis
$(N, D, C, b_{\text{train}}, b_{\text{deploy}})$ frontier is
incompletely mapped at frontier scale; we return to the
mixed-precision recipe in \cref{sec:mixed-precision}.

\subsection{Frontier-2026 scaling map}
\label{sec:scaling-map}

Tying the threads together: scaling at the 2026 frontier is a
four-axis surface $(C_{\text{train}}, C_{\text{test}}, b, q)$ where
$C_{\text{train}}$ is pre-training FLOPs, $C_{\text{test}}$ is
inference-time-compute budget, $b$ is training/deployment precision,
and $q$ is the quality-mixture composition of $D$ (covered separately
in \cref{sec:pre-training-data}; see Data Mixing Laws there). Frontier
labs jointly optimize all four. Where they agree: Chinchilla's
parametric form \cref{eq:cc-paraloss} fits, the $C \approx 6 N D$
formula is universal, \textmu P-inspired HP transfer is broadly
adopted (with disclosure asymmetry; see \cref{sec:scaling-mup}), \glsterm{fp8}{FP8}
or lower is the new training-precision floor where the stack supports
it. Where they diverge: tokens-per-parameter is now a
deployment-conditional choice in $[20, 2000]$ rather than a
universal $\approx 20$, and the $C_{\text{train}}$-vs-$C_{\text{test}}$
split is a per-product tuning that no closed-form law yet predicts.

\begin{contradiction}
Kaplan-vs-Chinchilla, on the methodological question, is settled:
the $0.73/0.27$ split was an LR-schedule artifact, the corrected
$\approx 0.50/0.50$ stands on three independent fits within bootstrap
of each other \citep[\S 3, Table 2]{hoffmann2022-chinchilla}. But
multiple ``optimal'' definitions are now in play. Pre-training-FLOP-
optimal (Chinchilla, $D / N \approx 20$); inference-cost-aware (Llama
3, $D / N$ in the hundreds to thousands; \citealp{meta-llama3});
inference-time-compute-aware (the o-series, R1, Gemini 2.5
\citep{gemini2-5}, Qwen 3 \citep{qwen3} with $C_{\text{test}}$
exposed); precision-aware (\citealp{scaling-laws-precision-2024} with
the over-training $\to$ harder-to-quantize trade-off pulling against
Llama-3-style $D/N \approx 1875$). The functional form
\cref{eq:cc-paraloss} survives; the prescription that comes out of
\cref{eq:cc-frontier} is now objective-conditional. ``Chinchilla is
wrong'' is a mis-framing; ``Chinchilla solves a different objective
than the one the industry actually optimizes'' is accurate. The
unresolved open question is the joint $(C_{\text{train}},
C_{\text{test}}, b, q)$ \glsterm{scaling-law}{scaling law} --- documented empirically, not
captured in closed form.
\end{contradiction}

\subsection{Forward link}
\label{sec:scaling-forward}

The scaling laws of this section dictate the budget split
$(N, D)$ at fixed $C$ and the HP-transfer protocol that makes any
chosen point trainable; they do not specify \emph{the optimizer} that
spends the FLOPs or the schedule that matches the LR-decay horizon
identified in \cref{sec:scaling-lrhorizon}. The next section
(\cref{sec:optimization}) walks the optimizer axis: \glsterm{adamw}{AdamW} as the
production default, the training-loss-conditional speedup claims of
\glsterm{lion}{Lion} and \glsterm{sophia}{Sophia} that did not hold on downstream evaluation, \glsterm{per-shape-rms-scaling}{Muon}'s
Newton-Schulz-orthogonalized matrix update as the only published
optimizer with $\sim 2\times$ compute-efficiency lift at \glsterm{llm}{LLM} scale,
and the WSD (\glsterm{wsd-schedule}{Warmup-Stable-Decay}) schedule that pairs with continual
pre-training and merging in \cref{sec:adaptation-and-merging}.
Optimizer choice and precision choice will turn out to be \emph{not}
orthogonal --- \glsterm{per-shape-rms-scaling}{Muon}'s update RMS exceeds BF16's high-precision range
without weight decay, exactly the optimizer$\times$precision coupling
\cref{sec:scaling-precision} flagged --- and the LR-schedule lesson
from \cref{sec:scaling-lrhorizon} reappears as the WSD cooldown phase
that loads quality data over the final tokens.
